\begin{textsample}{POS Dim 3 – ap – Score 13.00 – ap/ap\_inf\_13.txt}  \label{ex:f3_pos_148}
V – Campo de Experiência : Espaços, Tempos, Quantidades, Relações e Transformações As \textbf{crianças} vivem inseridas em espaços e tempos de diferentes dimensões, em um mundo constituído de fenômenos naturais e socioculturais.
Desde muito \textbf{pequenas}, elas procuram se situar em diversos espaços ( rua, bairro, cidade etc. ) e tempos ( dia e noite ; hoje, ontem e amanhã etc. ). \textbf{Demonstram} também \textbf{curiosidade} sobre o mundo físico ( seu próprio corpo, os fenômenos atmosféricos, os animais, as plantas, as transformações da natureza, os diferentes tipos de materiais e as possibilidades de sua manipulação etc. ) e o mundo sociocultural ( as relações de parentesco e sociais entre as pessoas que conhece ; como vivem e em que trabalham essas pessoas ; quais suas tradições e seus costumes ; a diversidade entre elas etc. ).
Além disso, nessas experiências e em muitas outras, as \textbf{crianças} também se deparam, \textbf{frequentemente}, com conhecimentos matemáticos ( contagem, ordenação, relações entre quantidades, dimensões, medidas, comparação de \textbf{pesos} e de comprimentos, avaliação de distâncias, reconhecimento de formas geométricas, conhecimento e reconhecimento de numerais cardinais e ordinais etc. ) que igualmente aguçam a \textbf{curiosidade}.
Portanto, a Educação \textbf{Infantil} precisa promover experiências nas quais as \textbf{crianças} possam fazer observações, \textbf{manipular} objetos, investigar e explorar seu entorno, levantar hipóteses e consultar fontes de informação para buscar respostas às suas \textbf{curiosidades} e indagações.
Assim, a instituição escolar está criando oportunidades para que as \textbf{crianças} ampliem seus conhecimentos do mundo físico e sociocultural e possam utilizá -los em seu cotidiano.
Direitos de aprendizagem com foco no campo de experiência espaços, tempos, quantidades, relações e transformações : - CONVIVER com \textbf{crianças} e \textbf{adultos} e com eles criar estratégias para investigar o mundo social e natural, \textbf{demonstrando} atitudes positivas em relação a situações que envolvam diversidade étnico-racial, ambiental, de gênero, de língua e de religião. - \textbf{BRINCAR} com materiais e objetos cotidianos, associados a diferentes papéis ou cenas sociais, e com elementos da natureza que apresentam diversidade de formas, \textbf{texturas}, cheiros, cores, tamanhos, \textbf{pesos}, densidades, experimentando possibilidades de transformação. - PARTICIPAR de atividades que oportunizem a observação de contextos diversos, atentando para características do ambiente e das histórias locais, utilizando ferramentas de conhecimento e instrumentos de \textbf{registro}, orientação e comunicação, como bússola, lanterna, lupa, máquina \textbf{fotográfica}, gravador, filmadora, projetor, computador, celular, dentre outros. - EXPLORAR e identificar as características do mundo natural e social, nomeando -as, reagrupando -as e ordenando -as, segundo critérios diversos. - EXPRESSAR suas observações, hipóteses e explicações sobre objetos, organismos vivos, fenômenos da natureza, características do ambiente, personagens e situações sociais, registrando -as por meio de desenhos, fotografias, gravações em áudio e vídeo, escritas e outras linguagens. - CONHECER -SE e construir sua identidade pessoal e cultural, identificando seus próprios interesses na relação com o mundo físico e social, apropriando -se dos costumes, das crenças e tradições de seus grupos de pertencimento e do patrimônio cultural, artístico, ambiental, científico e tecnológico.

% matched lemmas: adulto, brincar, criança, curiosidade, demonstrar, fotográfico, frequentemente, infantil, manipular, pequeno, peso, registro, textura
\end{textsample}
